\chapter{Introdução}
\label{cap:Introducao}

A geração, o tráfego, o armazenamento, o processamento, a transformação e o consumo de dados são características marcantes da comunicação e relações humanas. Sendo a espinha dorsal da economia digital (bem como extremamente relevante para a economia tradicional, dado o peso que a tecnologia e seu desenvolvimento têm sobre as atividades quer dos setores primário, secundário e terciário), os dados e a informação deles derivados são valiosos em seus diversos formatos. A despeito do peso da imagem -- seja estática ou em fluxo de vídeo -- e do som nesse universo contemporâneo que progressivamente busca o estímulo dos sentidos como forma de promoção de diversos valores e produtos, é inegável que dados textuais formam uma grande massa informacional.

Sejam em formato textual por natureza ou dados representados via texto, quer gerados por robôs ou produzidos por humanos, longos ou curtos, estruturados/tabulares, semi-estruturados ou não estruturados, esse volume de informação jamais disponível em qualquer tempo de nossa história humana tem o potencial de, ao ser analisado, servir como ferramenta útil em melhoramento de processos, tomada de decisão, difusão do conhecimento e outros infindos meios de enriquecimento de qualidade de oportunidades.

Não obstante, diante de uma torrente imensurável de dados, a capacidade de busca, análise, processamento e compreensão humana se amiúda por limitações diversas, sobretudo biológicas. Em contextos organizacionais, privados ou públicos, por exemplo, a geração de dados textuais que descrevem rotinas diversas e compõem processos, bem como o armazenamento da informação-fim, produto desses fluxos, acaba por compor um volume inestimável de informação que implica em prejuízo tanto em casos de sua perda quanto nos casos de falta de sua análise. A impossibilidade de se observar esse corpo de dados, em busca de informações, reduz o potencial de gerar conhecimento que viabilize mudanças e tomadas de decisão, rarefazendo ações latentes de aprimoramento da entidade em questão e da eficiência de seus processos internos.

Em um contexto como o da \ac{Alern}, são gerados e mantidos dados pertinentes a uma diversidade de escopos que são de interesse dos cidadãos, seja para fins de prestação de contas, habilitação ao exercício da cidadania ou oferecimento mais eficaz de serviços a serem consumidos pela população. Iniciativas como o Portal da Transparência Legislativa, no ar desde fevereiro de 2024 \cite{ALERN:3}, têm a preocupação de viabilizar o acesso do cidadão a parte dessa massa de dados que sejam de livre acesso (dados institucionais, receitas, despesas, servidores, contratos, demandas relativas às Leis de Acesso à Informação e de Responsabilidade Fiscal, pagamento de Diárias e planejamento e execução orçamentária).

É de relevância, contudo, apontar que a disponibilização desses dados, em específico, grandemente se facilita por sua natureza estruturada/semiestruturada, sendo também este exato aspecto que torna intuitiva e direta a forma pela qual os interessados interagem com as ferramentas utilizadas no seu acesso. Tal característica não é compartilhada por toda a gama dos dados produzidos pela \ac{Alern}, os quais são naturalmente não estruturados. Como parte do rol de dados não estruturados, por exemplo, está o que se poderia afirmar ser o produto fim da própria atividade legislativa (razão mor da existência da \ac{Alern}), a saber, o corpus de documentos composto pelas leis que regem toda a estrutura governamental no Estado do Rio Grande do Norte, a despeito de necessariamente seguirem o rigor característico de atos normativos. \cite{LeiRedacao:1}

Com efeito, o volume maior de informações produzidas pela \ac{Alern} é de natureza não estruturada, tanto as voltadas para o público externo quanto as de uso interno. A atividade legislativa, já mencionada, produz as leis estaduais, resoluções e outros atos normativos de aplicação em âmbitos diversos, atividade essa que, além do texto final da peça produzida, resulta em outras peças intermediárias durante a tramitação de proposições legislativas, como os textos das diversas proposições (projetos de lei, projetos de emenda constitucional, projeto de resolução e afins). As publicações diárias do Diário Oficial Eletrônico, igualmente, compõem um volume incremental e constante de novos documentos textuais, com diversas informações relativas a atividades gerais da \ac{Alern}, gerenciamento de recursos, editais e publicações em geral.

Paralelamente a isso, ainda no que se refere à atividade parlamentar, outra fonte de dados não estruturados de interesse público é o conjunto de transcrições das variadas interações dos Deputados em suas atuações. Os diversos tipos de Sessão Plenária (ordinárias, solenes, extraordinárias, dentre elas), as reuniões das Comissões, as Consultorias Públicas, as Audiências Públicas e demais ocasiões das reuniões parlamentares são registradas em forma de atas e têm suas gravações de vídeo transcritas automaticamente, tornando seu conteúdo disponível em formato de texto.

Os setores de tecnologia da Casa Legislativa atuam de forma notória constantemente para a disponibilização desse conteúdo ao público, incluindo busca textual por termos específicos e palavras chave. Não obstante, a consulta a documentos limitada pela ocorrência de palavras ou expressões tende a não incluir no conjunto de resultados documentos com conteúdo potencialmente relevante para a intenção do interessado, visto que podem não conter os termos ou expressões exatos utilizados como critérios de busca. Destarte, uma forma de recuperação de conteúdo que seja mais ampla e flexível poderia aprimorar grandemente as formas de disponibilização de conteúdo e garantir resultados mais satisfatórios ao público, colaborando para o cumprimento da obrigação de transparência do exercício parlamentar na \ac{Alern}.

De forma a atender a essa necessidade de consideração do contexto da consulta, do sentido para além das palavras utilizadas como parâmetros, uma possível abordagem é considerar semanticamente os termos de busca utilizados. Ao se utilizar a similaridade semântica entre os documentos sendo consultados e os dados inseridos pelo usuário como critérios de filtragem, amplia-se a cobertura de revocação, garantindo resultados de busca mais robustos e satisfatórios; isto partindo do pressuposto que a intencionalidade do consultante é em meramente filtrar dentre o grande volume de documentos um subconjunto específico dos tais que lhe seja de interesse.

Todavia, convém ponderar que a mera filtragem de dados, embora extremamente relevante no contexto em questão, providencia meramente a redução da dimensão do problema (que é o volume de dados em si) sem tratar de uma solução, no sentido mais estrito de sua definição. A interpretação dos dados resultantes dessa seleção feita requer uma segunda camada de intervenção, seja por parte do próprio interessado ou por meio de algum outro procedimento computacional. Ainda mais crítica se torna esse segundo processo nos casos em que as informações buscadas estejam espalhadas por diversos documentos do subcorpus resultante da filtragem, requerendo sua identificação e síntese. A combinação dessas informações ocasionaria a geração de informação útil que, embora seja oriunda do conjunto de dados preexistente, ao mesmo tempo o transcende, sendo mais do que a soma de suas partes. Tornar-se-ia, pois, uma tarefa, não somente de recuperação de dados, mas de geração de informações através de inferência a partir daqueles.

Os próprios termos utilizados para descrever a tarefa – geração de informação – acabam por apontar na direção de que meios poderiam ser utilizados para sua execução, a saber, a aplicação de modelos generativos. Através dessa abordagem, o subconjunto de dados obtido através de uma filtragem por meio de similaridade semântica seria utilizado por um \textit{\ac{LLM}} para, guiado pela intencionalidade de busca do interessado, produzir uma porção de texto que contenha o resultado desejado. Especificando ainda mais a solução, pode-se optar, como uma forma direta de se obter uma estrutura nesse formato, a agregação dos critérios de filtragem dos documentos inseridos pelo consultante com a intencionalidade em relação ao que se espera do resultado do processo como um todo, sendo uma possibilidade viável uma configuração de pergunta e resposta. A pergunta serve ao mesmo tempo como critério para escolha de documentos e como parâmetro fornecido ao \ac{LLM}, junto com os documentos selecionados, para a formulação de uma resposta, que seria a síntese direcionada do conteúdo de dados obtidos.


\section{Justificativa}

Diante da exposta profusão de dados textuais de interesse público gerados e mantidos pela \ac{Alern}, do interesse da instituição em sua disponibilidade aos cidadãos, assim como da atual inexistência de ferramentas que permitam consulta a dados com flexibilidade nos moldes expostos, torna-se evidente a necessidade premente de uma ferramenta de consulta de conteúdo textual adaptável aos diversos conjuntos de dados textuais da \ac{Alern}. De modo a alcançar esse objetivo, tal ferramenta, necessariamente, precisa de uma arquitetura que seja capaz de abarcar dados com volume e conteúdo dinâmicos, em contínuo incremento, o que implica em uma abordagem que não dependa de alteração/treinamento mediante alteração drástica no conteúdo do conjunto de dados, considerando as alterações constantes de documentos existentes e inserções diárias no corpus documental do órgão.

Ademais, dadas as evoluções crescentes nos sistemas computacionais e a adição de novas aplicações ao conjunto das já existentes, é recorrente a criação de novos conjuntos/coleções de dados textuais de forma, natureza e conteúdo completamente distintos do que já está estabelecido. Torna-se, portanto, a flexibilidade ainda mais necessária, de forma a evitar a necessidade de redesenho constante da base comum de operações compartilhadas entre os diversos conjuntos de dados.

\section{Objetivos}
\label{sec:Introducao.1}

O conceito definidor da solução proposta consiste em um sistema de perguntas e respostas (comumente designado como \textit{Q\&A System}), nos moldes de um \textit{chatbot}, baseado em \textit{Retrieval-Augmented Generation} (RAG), que se utiliza de \textit{Large Language Models} (LLMs) alimentados com documentos textuais, com o intuito de gerar respostas a perguntas dentro de um domínio coberto por um corpus de documentos relevantes.

Por meio do estabelecimento de uma estrutura base que recebe \textit{prompts} com uma pergunta e o conteúdo de um conjunto de documentos selecionados, se torna possível a utilização de uma mesma arquitetura utilizável para consultas em diversas bases de conhecimento distintas, podendo cobrir tantos domínios de conteúdo textual quantos forem achados relevantes. Com isso, por meio de se criar, para cada conjunto de dados textuais de interesse, uma coleção de documentos a ser armazenada em uma \textit{Vector Store} (um banco de dados otimizado para uso com vetores de \textit{embeddings}), a mencionada estrutura pode ser utilizada de forma generalista, alterando somente a seleção de documentos e a pergunta a serem processadas pelo modelo de linguagem. 

Com essa estratégia, evitada fica a necessidade de \textit{fine-tuning} de algum modelo textual para se adaptar ao conteúdo nativo do conjunto de dados da \ac{Alern}, garantindo resultados positivos mesmo quando da alteração de suas instâncias ou da inclusão de novos dados com teor completamente distinto do já existente. A cada novo documento inserido nas coleções já existentes, bem como a cada nova coleção acrescentada ao conjunto de dados de interesse, basta fazer a preparação dos dados, garantindo igualmente a sua inclusão no banco de vetores, e a estrutura de consulta, recuperação de documentos relevantes e geração de conteúdo direcionado está pronta para atuar de forma a incluir as informações novas no sistema.

\subsection{Objetivo Geral}
\label{subsec:Introducao.1.1}
Construir de uma ferramenta de perguntas e respostas no formato de um \textit{chatbot}, com o intuito de buscar conteúdo e gerar respostas sintetizadas/sintéticas, que abarque de forma flexível diversos conjuntos de dados textuais de interesse público e interno no contexto da \ac{Alern}.

\subsection{Objetivos Específicos}
\label{subsec:Introducao.1.2}

%AFAZER: Alterar os textos dos objetivos para apontarem mais para os "artefatos", os "produtos" de cada uma das fases do trabalho

Preparar uma base de conhecimento a partir de um conjunto de documentos de um domínio específico da \ac{Alern};

Definir uma estrutura de recuperação de fragmentos desses documentos com base em uma pergunta fornecida pelo usuário;

Desenhar uma função de estimativa de relevância dos documentos;

Implementar uma solução de geração/síntese do conteúdo recuperado com foco em responder a pergunta;

Integrar os elementos desenvolvidos, de forma coesa, incluindo uma interface de usuário para a realização da consulta; 

Realizar testes e avaliar o desempenho e a precisão das informações recuperadas e geradas pelo \textit{chatbot}; 

Ajustar todos os elementos da estrutura de forma a minimizar problemas encontrados nos testes e avaliações.